\chapter{Results and Analysis}\label{results}
This chapter is dedicated to presenting the results of our experiments. We will begin by showcasing the outcomes of the hyperparameter tuning process for the evaluation models and the augmentation techniques. Then we will present the results of both the classification and detection tasks for the basic augmentations. Finally, we will dive into the results of using GAN to enhance the performance of edge cases. To improve the visual representation, a smoothing average window of 0.7 was applied to all graphics.

\section{Baseline and Hyperparameter tuning}
In this section, we present the baseline for the YOLOv5s and EfficientNet for the given task, along with the optimal hyperparameters for the augmentation techniques. The baseline experiments were conducted without any data augmentation and involved training the models both with and without pre-trained weights. The evaluation metrics used to measure the performance of the models include F1 score Top-3 and mAP@[0.5:0.95].
\subsection{Pretrain weights and learning rate} \label{results:pretrained}
We perform learning rate hyperparameter tuning for our evaluation models and consider using either random weights or pretrained weights. Selecting the best possible evaluation model is crucial as it can greatly affect the evaluation of our augmentation techniques. An underperforming model may produce incorrect predictions, leading to inaccurate results and unreliable conclusions.

\subsubsection{YOLOv5s}
Two YOLOv5 models were trained without augmentation for 100 epochs: one with pretrained weights from the COCO dataset, and one without pretrained weights. Results in Figure \ref{fig:baseline_yolo}a) demonstrate that the pretrained model outperforms the non-pretrained model. This suggests that the general features learned during pretraining, such as edges, corners, and textures, are still useful for detection in our dataset even though the COCO is not specific to aerial imagery \cite{cocodataset}.

To further optimize performance, a hyperparameter tuning experiment was conducted to determine the optimal learning rate. The results, shown in Figure \ref{fig:baseline_yolo}b), indicate that the best learning rate was 0.01, which achieved an F1-Score top-3 of 0.1773. To summarize, we have determined that YOLOv5s with pretrained weights and a learning rate of 0.01 is the baseline for the deteciton task.

\begin{figure}[!h]
\centerline{\subfigure[With and without pretrained weights]{\includegraphics[width=12cm]{src/Results/baseline/base_detec_pre.png}}}
\centerline{\subfigure[Learning rate tuning]{\includegraphics[width=12cm]{src/Results/baseline/detec_lr.png}}}
\caption{Illustration of (a) YOLOv5s with and without pretrained weights and (b) on different learning rate}
\label{fig:baseline_yolo}
\end{figure}

% \begin{figure}[!h]
% 	 \centerline{\includegraphics[width=0.7\textwidth]{src/Results/baseline/base_detec_pre.png}}
% 	 {\caption{mean average precision of Yolov5 Lite with and without pretrained weights}\label{fig:base_detec_pre}}
% \end{figure}

% \begin{figure}[!h]
% 	 \centerline{\includegraphics[width=0.7\textwidth]{src/Results/baseline/detec_lr.png}}
% 	 {\caption{Learning Rate Hyperparatmeter tuning for Yolov5}\label{fig:detec_lr}}
% \end{figure}
\subsubsection{EfficientNet}

The EfficientNet was trained for 200 epochs using both random-initialized weights and pretrained weights from the Imagenet dataset, without any augmentation. The results, as illustrated in figure \ref{fig:baseline_effi}a), demonstrate that the pretrained model outperformed the model without pretrained weights, suggesting that learning transferability exists even when the datasets lack significant common features.

After conducting hyperparameter tuning within a logarithmic search space between $10^{-4}$ and $10^{-8}$, we determined that the optimal learning rate for EfficientNet version B0 with pretrained weights was $10^{-5}$, as demonstrated in figure \ref{fig:baseline_effi}b). This yielded an F1-Score top-3 of 0.3434. Thus, we established the baseline for the classification task as using EfficientNet version B0 with pretrained weights and a learning rate of $10^{-5}$.

% \begin{figure}[!h]
% 	 \centerline{\includegraphics[width=1\textwidth]{src/Results/baseline/base-class-pre.png}}
% 	 {\caption{F1-Score top-3 of EfficienNet B0 with and without pretrained weights}\label{fig:base-class-pre}}
% \end{figure}

% \begin{figure}[!h]
% 	 \centerline{\includegraphics[width=1\textwidth]{src/Results/baseline/class_lr.png}}
% 	 {\caption{Learning Rate Hyperparatmeter tuning for EfficientNet}\label{fig:class_lr}}
% \end{figure}

\begin{figure}[!ht]
\centerline{\subfigure[With and without pretrained weights]{\includegraphics[width=12cm]{src/Results/baseline/base-class-pre.png}}}
\centerline{\subfigure[Learning rate tuning]{\includegraphics[width=12cm]{src/Results/baseline/class_lr.png}}}
\caption{Illustration of EfficientNet version B0 (a) with and without pretrained weights and (b) on different learning rates}
\label{fig:baseline_effi}
\end{figure}

\subsection{Augmentation Hyperparameter Tuning}
The augmentation techniques that required tuning included CLAHE, Color Jittering, Griddropout, Cutout, Sharpen, Motion Blur, Perspective, and RandAugment. We conducted the tuning process using grid search and random grid methods. Table \ref{tab:multicol} presents a summary of the optimal hyperparameters obtained from the tuning process, and a more detailed description of the hyperparameter search can be found in the \hyperref[app:Appendix]{Appendix}.

\begin{table}[!ht]
\begin{center}
\caption{Hyperparameter tuning results for augmentations} \label{tab:multicol}
\begin{tabular}{|c|c|c|c|}
    \hline
    \textbf{Method} & \textbf{Hyperparameter} & \textbf{\makecell{Optimal\\ values}} & \textbf{Search Space} \\ \hline
    \multirow{2}{*}{\hyperref[tab:clahe_hp]{CLAHE}}& Clipping Limit & 16 & [2, 4, 8, 16, 32]\\ \cline{2-4} 
    & Filter Size & 4 & [2, 4, 8, 16, 32]\\ 
    \hline
    \multirow{4}{*}{\hyperref[tab:colorjittering_hp]{Color Jittering}}& Brightness limit& 0.8 & \makecell{Randomly between \\ 0.2 and 0.8}\\ \cline{2-4} 
    &Hue & 0.42 & \makecell{Randomly between \\ 0.2 and 0.8}\\ \cline{2-4} 
    &Contrast limit & 0.54 & \makecell{Randomly between \\ 0.2 and 0.8}\\\cline{2-4} 
    &Saturation limit & 0.65 & \makecell{Randomly between \\ 0.2 and 0.8}\\ \hline
    \hyperref[tab:griddropout_hp]{GridDropout}  & Hole Ratio & 0.5 & [0.2, 0.3, 0.4, 0.5, 0.6]\\ \hline
    \multirow{2}{*}{\hyperref[tab:sharpen_hp]{Sharpen}}& Alpha range & (0.6, 0.8)  & \makecell{0.2-interval between 0.2 and 0.8 \\ with window sliding step of 0.1}\\ \cline{2-4} 
    & Lightness range & (0.6, 1.0)  & \makecell{0.2-interval between 0.2 and 0.8 \\ with window sliding step of 0.1}\\ \hline
    \hyperref[tab:motion_blur_hp]{Motion blur} & Blur Limit & 21 & [11, 21, 31, 41, 51, 61, 71] \\ \hline
    \hyperref[tab:perspective_hp]{Perspective} & Scale limit & 0.2 & [0.2, 0.3, 0.4, 0.5, 0.6, 0.7] \\ \hline
    RandAug & \makecell{Number of \\combinations} & 2 & [2, 3, 4] \\ \hline
\end{tabular}
\end{center}


\end{table}


\section{Results of basic augmentation techniques}
In this section, we present the outcomes of applying fundamental augmentations, which include Geometric, Color-Space, Cutout, and RandAugment techniques. The section is divided into two parts, with the first one showcasing the results of the detection task using YOLOv5s, and the second one presenting the findings of the classification task using EfficientNet B0.
\subsection{Detection Task}
We trained the pretrained YOLOv5 model using a learning rate of 0.01 for a total of 100 epochs. Table \ref{tab:detection-result} provides a summary of the resulting performance metrics, including the average mean precision at 0.5 IoU threshold (mAP@0.5), average mean precision at IoU threshold between 0.5 and 0.95 (mAP@[0.5:0.95]), best precision, and best recall. Following the overall results, we proceed to analyze the outcomes of each group of augmentations, namely Geometric Augmentation, Color-Space Augmentation, Cutout, and RandAugment.

\subsubsection{Overall}
To summarize, most of the geometric augmentations negatively impacted the model's performance, with the exception of CenterCrop. CenterCrop achieved the highest mAP@0.5 and mAP@0.5:0.95, respectively with values of 0.3436 and 0.5214, which is nearly two times better than the baseline performance. On the other hand, color-space augmentations generally improved the model's performance, with the best augmentation technique being ToSepia, which scores 24\%  in mAP@[0.5:0.95], a notable 6.29\% improvement over the baseline. Cutout and GridDropout had a notable negative impact on the model's performance, resulting in a decrease of 2\% and 7.5\%, respectively. However, the worst-performing augmentation technique was RandAugment, which produced a staggering 10\% decrease in the mAP.

\begin{table}[!htb]
    \centering
    \caption{Results of augmentations on YOLOv5s at fine-tuned learning rate $10^{-3}$}    \label{tab:detection-result}
    \begin{tabular}{|l|llll|}
    \hline
        \textbf{Method} & \textbf{\makecell{mAP@[0.5:0.95]}} & \textbf{\makecell{mAP@0.5}} & \textbf{\makecell{Precision}} & \textbf{\makecell{Recall}} \\ \hline
        Baseline & 0.1773 & 0.2879 & 0.9298 & 0.6391\\ 
        \hline
        CenterCrop 0.2 & \textbf{0.3436} & \textbf{0.5214} & 0.7719 & 0.6865 \\ 
        Flipping & 0.1584 & 0.2772 & \textbf{0.9576} & \textbf{0.6931} \\ 
        Motion Blur & 0.1517 & 0.2256 & 0.4693 & 0.2177 \\ 
        Perspective & 0.1244 & 0.1999 & 0.2695 & 0.2256 \\
        Rotation & 0.1427 & 0.2303 & 0.3325 & 0.2665 \\ 
        \hline
        CLAHE & 0.2152 & 0.3369 & 0.7346 & 0.4768 \\ 
        ChannelShuffle & 0.235 & 0.3619 & 0.7651 & 0.575 \\ 
        GaussNoise & 0.2077 & 0.3236 & 0.7858 & 0.525 \\ 
        Solarize & 0.1612 & 0.2683 & 0.3872 & 0.273 \\
        Sharpen & 0.2402 & 0.3604 & 0.7413 & 0.4545 \\ 
        ToGray & 0.2394 & \textbf{0.363} & \textbf{0.7741} & \textbf{0.5432} \\ 
        ToSepia & \textbf{0.246} & 0.3757 & 0.7404 & 0.5437 \\ 
        \hline
        Cutout & \textbf{0.1503} & \textbf{0.2481} & 0.4137 & 0.2427 \\ 
        GridDropout & 0.1023 & 0.1535 & 0.3643 & 0.1549 \\ 
        \hline
        RandAugment & 0.07167 & 0.1347& 0.5115& 0.1305 \\
    \hline
    \end{tabular}
\end{table}

\subsubsection{Geometric Augmentation}\label{result:dec:geo}
The mAP@[0.5:0.95] metric for the geometric augmentation technique over 100 epochs is visually represented in Figure \ref{fig:yolo_geo}. As demonstrated in the results, CenterCrop augmentation technique outperforms all other augmentation techniques. This observation confirms our initial assumption that an augmentation technique that ensures center invariance can considerably enhance the model's performance in terms of both bounding box detection and class prediction, as depicted in Figure \ref{fig:geo_loss}.

\begin{figure}[!h]
	 \centerline{\includegraphics[width=0.9\textwidth]{src/Results/detection/yolo_geo.png}}
	 {\caption{Mean average precision in threshold interval between 0.5 and 0.95 of YOLOv5s with geometric augmentations}\label{fig:yolo_geo}}
\end{figure}

Although it is reasonable to assume that motion blur would not improve a model's performance due to the added noise, the fact that the Rotation, Flipping, and Perspective augmentations decreased the overall performance of the model is counter-intuitive. These transformations are generally expected to increase the robustness of the training set, but they may actually have an adverse effect on the model's ability to learn.

One possible explanation for this unexpected outcome is that YOLOv5 is a bounding box detection model. A modification in the geometric shape of the object can result in a more varied squared shape of the bounding box, which can slow down the training convergence. This assumption is supported by the results shown in figure \ref{fig:geo_loss}, where we observe that although the classification loss did not substantially increase with the augmentations (except for Perspective), the box loss of Flipping, Perspective, and Rotation decreased much slower than the baseline. Hence, the architecture of the YOLOv5 model needs to be considered while selecting augmentation techniques for improving the model's performance.

\begin{figure}
\hfill
\subfigure[Box loss]{\includegraphics[width=7cm]{src/Results/detection/geo_box_loss.png}}
\hfill
\subfigure[Classification loss]{\includegraphics[width=7cm]{src/Results/detection/geo_class_loss.png}}
\hfill
\caption{Illustration of (a) box loss and (b) classification loss in geometric augmentation group}
\label{fig:geo_loss}
\end{figure}

\subsubsection{Color-Space Augmentation}

\begin{figure}[!h]
	 \centerline{\includegraphics[width=0.9\textwidth]{src/Results/detection/yolo_color.png}}
	 {\caption{Mean average precision of YOLOv5s with Color-Space Augmentations}\label{fig:yolo_color}}
\end{figure}

Figure \ref{fig:yolo_color} provides a visualization of the mAP@[0.5:0.95] metric for color-space augmentations. In general, all color-space augmentation techniques significantly improve the performance of YOLOv5.

The plot in Figure \ref{fig:color_loss} illustrates the trade-off between classification loss and box-loss in color-space augmentation. The extra sharpness and contrast introduced by this augmentation can significantly improve the model's ability to predict the correct class label. However, this comes at a cost of decreased accuracy in detecting the bounding box. This indicates that the model with color-augmentation could also learn features of the background to aid in object classification.

\begin{figure}[!h]
\hfill
\subfigure[Box loss]{\includegraphics[width=7cm]{src/Results/detection/color_box_loss.png}}
\hfill
\subfigure[Classification loss]{\includegraphics[width=7cm]{src/Results/detection/color_class_loss.png}}
\hfill
\caption{Illustration of (a) box loss and (b) classification loss in Color-Space Augmentations}
\label{fig:color_loss}
\end{figure}

While CLAHE and Sharpen work by enhancing the contrast and brightness of images, making it easier for the model to detect and classify objects, the fact that ToSepia and ToGray augmentations also improve performance is noteworthy. Interestingly, ToSepia outperforms all other color-space augmentations, indicating that reducing color information can enhance the model's detection ability. By reducing the variances in the three channels and consolidating them into one dimension, the model can focus more on the object's shape rather than color, resulting in improved performance.

\subsubsection{Cutout and RandAugment}
Figure \ref{fig:yolo_cutout} shows that the performance of the model is reduced by both Cutout and GridDropout. Furthermore, RandAugment causes the most significant reduction in performance, with the largest margin observed among all augmentation techniques.

\begin{figure}[!h]
	 \centerline{\includegraphics[width=0.9\textwidth]{src/Results/detection/yolo_cutout.png}}
	 {\caption{Mean average precision of YOLOv5s with Cutout Augmentations and RandAugment}\label{fig:yolo_cutout}}
\end{figure}

The reasons for the decrease in performance caused by Cutout and Dropout can be relatively intuitive, as they are non-label-preserving augmentation techniques that may remove crucial information required for the model to detect and classify objects accurately. Figure \ref{fig:cutout_loss} demonstrates that Cutout produces a significant increase in the box loss, even more so than GridDropout, which actually removes more information. Meanwhile, RandAugment performs poorly in both classification and box losses, introducing too many inconsistencies into the model and hindering its ability to identify crucial features.

\begin{figure}[!h]
\hfill
\subfigure[Box loss]{\includegraphics[width=7cm]{src/Results/detection/cutout_box_loss.png}}
\hfill
\subfigure[Classification loss]{\includegraphics[width=7cm]{src/Results/detection/cutout_class_loss.png}}
\hfill
\caption{Illustration of (a) box loss and (b) classification loss in Cutout augmentations and RandAugment}
\label{fig:cutout_loss}
\end{figure}
\subsection{Classification Task}
Table \ref{tab:classification-result} presented Top-3 and Top-1 metrics obtained from the classification outcomes, which are categorized into four sections, namely baseline, geometric transformation, color-space augmentation, and cutout-related augmentation. The training process was carried out over 200 epochs and implemented a learning rate of 0.0001. The highest metric in each group is emphasized in bold. Following the overall results, we proceed to analyze the outcomes of each group of augmentations, namely Geometric Augmentation, Color-Space Augmentation, Cutout, and RandAugment.

\begin{table}[!h]
\caption{Result of augmentation on EfficinetNet B0 at fine-tuned learning rate $10^{-5}$}    \label{tab:classification-result}

    \centering
    \small
    \tabcolsep=0.11cm
    \begin{tabular}{|l|lll|lll|}
    \hline
    & \multicolumn{3}{c}{\textbf{Top-3 Metrics}} & \multicolumn{3}{c}{\textbf{Top-1 Metrics}} \\ \hline
        \textbf{Method} & \textbf{\makecell{F1-score \\Top-3}} & \textbf{\makecell{Accuracy \\ Top-3}} & \textbf{\makecell{Precision \\ Top-3}} & \textbf{\makecell{F1-score \\ Top-1}} & \textbf{\makecell{Accuracy \\ Top-1}} & \textbf{\makecell{Precision \\ Top-1}} \\ \hline
        Baseline & 0.3434 & 0.4977 & 0.3231 & 0.2784 & 0.3339 & 0.2264 \\ \hline
        Rotation & \textbf{0.4423} & \textbf{0.5876} & \textbf{0.432} & \textbf{0.3364} & \textbf{0.4586} & \textbf{0.2585} \\ 
        Flipping & 0.4208 & 0.5826 & 0.4104 & 0.3296 & 0.4351 & 0.2563 \\ 
        Perspective & 0.3577 & 0.5209 & 0.3339 & 0.286 & 0.3414 & 0.218 \\ 
        CenterCrop 0.2 & 0.3252 & 0.5111 & 0.3028 & 0.2829 & 0.3284 & 0.2204 \\ 
        CenterCrop 0.4 & 0.316 & 0.5205 & 0.3023 & 0.2807 & 0.3333 & 0.209 \\ 
        MotionBlur & 0.3347 & 0.5209 & 0.3233 & 0.3054 & 0.3486 & 0.2353 \\ \hline
        CLAHE & 0.3027 & 0.4654 & 0.2932 & 0.2569 & 0.3093 & 0.1971 \\ 
        Sharpen & 0.3393 & 0.5095 & 0.3168 & \textbf{0.306} & 0.3525 & 0.2486 \\ 
        ChannelShuffle & 0.314 & 0.4682 & 0.2945 & 0.2646 & 0.3188 & 0.2011 \\ 
        Solarize & 0.2354 & 0.4165 & 0.2173 & 0.2251 & 0.2191 & 0.1666 \\ 
        Gaussian Noise & \textbf{0.3498} & 0.5154 & \textbf{0.3289} & 0.2868 & \textbf{0.3546} & 0.2168 \\ 
        MotionBlur & 0.3347 & \textbf{0.5209} & 0.3233 & 0.3054 & 0.3486 & \textbf{0.2353} \\ 
        % Normalize & 0.3315 & 0.5123 & 0.3119 & 0.2778 & 0.3395 & 0.2152 \\ 
        ToGray & 0.2941 & 0.4637 & 0.2787 & 0.2638 & 0.3011 & 0.211 \\ 
        ToSepia & 0.329 & 0.4708 & 0.3183 & 0.2695 & 0.3436 & 0.2111 \\ \hline
        Cutout & 0.2769 & 0.4891 & 0.2167 & 0.3387 & 0.3199 & 0.3396 \\ 
        GridDropout & 0.2827 & 0.4894 & 0.2167 & 0.333 & 0.3154 & 0.3454 \\ \hline
        RandAugment & 0.2929 & 0.5243 & 0.3789 & 0.3598 & 0.3731 & 0.2174 \\ 
    \hline
    \end{tabular}
    
\end{table}

\subsubsection{Overall}
In summary, most of the augmentation techniques did not lead to an improvement in the performance of EfficientNet. However, Rotation, Flipping, Perspective, and Gaussian Noise augmentations did show positive impacts on the model's performance. In particular, Rotation outperformed all other augmentations in terms of metrics and improved the baseline F1-Score Top-3 by almost 10\%. While Gaussian Noise was the only augmentation in the color-space group that had a positive impact on the performance, the difference was relatively small, with an increase of 0.0064 in F1-Score Top-3. Finally, regularization techniques such as cutout and dropout worsened the model's performance.

\subsubsection{Geometric Augmentations}
Figure \ref{fig:class-geo} depicts the classification performance of EfficientNet B0 when using geometric augmentations. Interestingly, contrary to the findings in the \hyperref[result:dec:geo]{detection model}and our initial hypothesis, CenterCrop augmentation does not seem to improve the model's classification performance. This observation implies that the model may require contextual information from the background to accurately predict the class. For example, in the case of a flock of birds, it may indicate a specific bird class that is commonly seen in a group, such as Graugans. Additionally, the geographical background may appear more frequently in certain bird classes, such as Brandgans, which could also assist in accurate classification.

\begin{figure}[!h]
	 \centerline{\includegraphics[width=0.9\textwidth]{src/Results/classification/class-geo.png}}
	 {\caption{Performance of Efficientnet B0 with Geometric Augmentations with fined-tuned learning rate $10^{-5}$.}\label{fig:class-geo}}
\end{figure}

Alternatively, it is possible that EfficientNet is already designed to learn features at different scales and resolutions, and by reducing the information, the model may lose some crucial details for a specific class. Therefore, CenterCrop augmentation may not be as effective in improving the model's classification performance.

However, Rotation, Flipping, and Perspective augmentations demonstrate their effectiveness in helping the model learn features in different geometric positions. With rotation the the F1-score improved substantially by 0.1. 


\subsubsection{Color-Space Augmentations}

\begin{figure}[!h]
	 \centerline{\includegraphics[width=0.9\textwidth]{src/Results/classification/class-color.png}}
	 {\caption{Performance of Efficientnet B0 with Color-space Augmentations with fined-tuned learning rate $10^{-5}$.}\label{fig:class-color}}
\end{figure}
Figure \ref{fig:class-color} displays the outcomes of the impact of color-space transformation on the performance of EfficientNets. The overall effect of this type of transformation is either negligible or insignificant with the exception of Gaussian Noise, which has a minimal positive effect. A possible explanation for this observation is that the architecture of EfficientNet is inherently designed to be resilient to color variations. The model employs a feature extractor network that learns to extract meaningful features from the input image, which is then coupled with a scaling method to ensure proper normalization of the input data. This design choice enhances the model's invariance to changes in color-space, rendering color-space augmentation less effective in improving the model's performance.


\subsubsection{Cutout Augmentations and RandAugment}


The results presented in Figure \ref{fig:class-cutout} demonstrate the impact of the cutout transformations and RandAugment on the performance of EfficientNets. It can be observed that the cutout transformation causes a decrease in the network's performance due to the loss of information. Similarly, incorporating RandAugment in the training pipeline results in a decrease in the network's performance. This could be due to the extreme nature of the augmentation techniques used in RandAugment, which can potentially slow down the convergence of the network during training.

\begin{figure}[!h]
	 \centerline{\includegraphics[width=0.9\textwidth]{src/Results/classification/f1-score-other.png}}
	 {\caption{Performance of Efficientnet B0 with Cutout, GridDropout and RandAugment with fined-tuned learning rate $10^{-5}$.}\label{fig:class-cutout}}
\end{figure}









\section{GAN-generated augmentations}
In this section, we will examine the outcomes of the GAN-generated methods, with a particular focus on evaluating the image quality, diversity, and effectiveness of the generated images on edge cases. The primary objective of this analysis is to highlight the common issues that can arise when using GANs as an augmentation technique and to compare the performance of the three architectures with one another.

Because of their complex architecture and the extensive number of parameters that must be adjusted, Generative Adversarial Networks (GANs) often necessitate a significant amount of time to train. Additionally, the training process can be computationally intensive, requiring high-end hardware and substantial computational resources. Table \ref{tab:gan_comparison} provides an overview of the number of trainable parameters, training time, and epochs required for each GAN employed in this thesis.

\begin{table}[h]
\centering
\begin{tabular}{|c|c|c|c|}
\hline
\textbf{Model} & \textbf{Number of trainable parameters} & \textbf{Epochs} & \textbf{Training Time} \\
\hline
Vanilla GAN & 405574756 & 10000 & 6h 58m 24s \\
WGAN & 743760999 & 5000  & 1d 5h 31m \\
WGAN-GP & 7645523 & 3000 & 10h 4m 50s \\
\hline
\end{tabular}
\caption{Comparison of GAN models}
\label{tab:gan_comparison}
\end{table}

\subsubsection{Image quality} 
\begin{figure}[!h]
	 \centerline{\includegraphics[width=1\textwidth]{src/Results/gan/samples_gan.pdf}}
	 {\caption{Samples of generated images from the conditional GANs}\label{fig:gan_samples}}
\end{figure}
The samples of generated images from three different types of GANs are displayed in Figure \ref{fig:gan_samples}. Each row represents an example of an edge class, such as Schnatterente, Graugans, Wasserlinie mit Großalgen, and Schmarotzerraubmöwe. The first four columns show the input images used to train the different GAN models. 

The images generated from Vanilla GAN appear to be the most realistic compared to other types of GANs. However, it is also observed that Vanilla GAN is prone to getting trapped in a local optima, resulting in generated images that closely resemble a specific real image, as seen in the case of the Schnatterente class. 

As illustrated with small seeming random patches in WGAN, when trained on a small dataset, WGAN did not have enough information to learn the underlying distribution of the data and therefore generate unrecognizable images. This is in contrast to vanilla GAN, which may perform better on smaller datasets but may suffer from mode collapse. 

Wasserstein GAN with gradient penalty has been demonstrated to produce generated samples that are more similar to the distribution of the training data in comparison to Vanilla WGAN. Although the generated images still exhibit a deficiency in high-resolution details, they share some similarities with the objects in the dataset. For instance, the class of Schnatterente features white dots and yellow patches, while the class of Wasserlinie mit Großalgen displays a similar shade of blue as shown in \ref{fig:gan_samples}. This observation highlights the issue of hypersensitivity to hyperparameters in WGAN. In the absence of hyperparameter tuning for the model, it failed to converge rapidly.

\subsubsection{Diversity}
Figure \ref{fig:gan_divers} displays different samples generated by 3 types of GANs with varying noise inputs across 3 classes. It is evident that Vanilla GAN suffers from the mode collapse problem, as it only produces one repeated image with slight variations in brightness and hue. On the other hand, WGAN and WGAN with GP successfully generate diverse samples.

\begin{figure}[!h]
	 \centerline{\includegraphics[width=1\textwidth]{src/Results/gan/divers_gan.pdf}}
	 {\caption{Variety of generated pictures in three GANs}\label{fig:gan_divers}}
\end{figure}

\subsubsection{Performance} 

\begin{figure}[!h]
\hfill
\subfigure[F1-Score Top-3 for all classes \label{fig:overall_gan}]{\includegraphics[width=7cm]{src/Results/gan/overall_gan.png}}
\hfill
\subfigure[F1-Score Top-3 for only edge classes \label{fig:edge_gan}]{\includegraphics[width=7cm]{src/Results/gan/edge_gan.png}}
\hfill
\caption{F1-Score (a) for all classes and (b) for edge classes}
\label{fig:gan_performance}
\end{figure}

% \begin{figure}[!h]
% 	 \centerline{\includegraphics[width=0.7\textwidth]}
% 	 {\caption{}}
% \end{figure}
Figure \ref{fig:overall_gan} depicts the F1-Score of EfficientNet when trained with additional GAN-generated data, while Figure \ref{fig:edge_gan} displays the F1-Score for edge cases. Overall, the overall classification performance decreases with the use of GAN-generated data. One possible explanation is that the GAN-generated images may contain noise or inconsistencies that are absent in the original dataset, thereby confusing the model and making it more difficult to generalize to new examples. Additionally, GAN-generated images may not cover the entire distribution of the original dataset, leading to overfitting and poor generalization performance.



WGAN with GP causes the largest decrease in F1-Score Top-3, with a decrease of nearly 5\%. On the other hand, using data from Vanilla GAN leads to the smallest decrease in performance, with a decrease of only 1.8\%. For the edge classes, both WGAN and Vanilla GAN show an improvement in F1-Score Top-3. WGAN results in a larger improvement of 0.013, while Vanilla GAN yields a smaller improvement of 0.0014. However, it is important to consider the overall decrease in performance when using GAN-generated data. The primary reason for the improvement in edge cases may be the introduction of inconsistencies and new features, which the network then learns to recognize. However, this improvement may not outweigh the decrease in performance seen in the rest of the dataset. In the case of WGAN with GP, the overall decrease in performance is 0.0355, which may indicate that the network is overfitting to the GAN-generated data. In the case of Vanilla GAN, the overall decrease in performance is 0.0205, which is still significant but less than that of WGAN with GP.





